\documentclass[10pt,letterpaper]{article}

\usepackage[margin=0.72in]{geometry}
\usepackage[hidelinks]{hyperref}
\usepackage{enumitem}
\usepackage{titlesec}
\usepackage{xcolor}
\usepackage{tabularx}

\definecolor{accent}{HTML}{0F766E}
\definecolor{textgray}{HTML}{4B5563}

\pagestyle{empty}
\setlength{\parindent}{0pt}
\setlength{\parskip}{0pt}
\setlist[itemize]{leftmargin=1.2em, itemsep=2pt, topsep=3pt}

\titleformat{\section}
  {\large\bfseries\color{accent}}
  {}
  {0pt}
  {}
  [\vspace{2pt}\titlerule]
\titlespacing*{\section}{0pt}{12pt}{7pt}

\newcommand{\cvitem}[2]{
  \begin{tabularx}{\textwidth}{@{}X r@{}}
    #1 & {\color{textgray}#2}
  \end{tabularx}
}

\newcommand{\paper}[5]{
  \textbf{#1} \hfill {\color{textgray}#2}\\
  {\itshape #3}\\
  #4
  \begin{itemize}
    #5
  \end{itemize}
}

\begin{document}

\begin{center}
  {\LARGE \textbf{Yuting Yan}}\\
  \vspace{4pt}
  M.S. Student in Software Engineering, University of Science and Technology of China\\
  \vspace{4pt}
  \href{mailto:ytyan@mail.ustc.edu.cn}{ytyan@mail.ustc.edu.cn}
  \quad | \quad
  \href{https://ytyan-bot.github.io/ytyan.github.io/}{Personal Homepage}
\end{center}

\section{Research Interests}
Large language models and intelligent agents; brain-computer interfaces and EEG analysis; reliable medical AI under cross-subject and cross-domain distribution shifts; test-time adaptation; multimodal clinical signal analysis.

\section{Education}
\cvitem{\textbf{University of Science and Technology of China}\\
M.S. in Software Engineering}{2024--2027}
\vspace{6pt}

\cvitem{\textbf{Southwest Jiaotong University}\\
B.E. in Artificial Intelligence}{2019--2023}

\section{Selected Publications}
\paper
  {Metric-Aware Test-Time Adaptation for Cross-Subject Multimodal Epileptiform-Discharge Detection}
  {ACM MM 2026}
  {First Author}
  {A label-free test-time adaptation framework for unseen-patient epileptiform-discharge detection. JMTR adapts candidate rankings for AUPRC, while DGC calibrates subtype posteriors under patient-dependent prior shift.}
  {
    \item NeuroMM-2026 Challenge: 3rd place in binary spike detection.
    \item NeuroMM-2026 Challenge: 3rd place in five-class subtype classification.
    \item Frozen backbone, no target labels, post-hoc adaptation.
  }

\paper
  {Reparameterizing Mamba via Frequency-Induced Topological Conduction for Medical Image Segmentation under Clinical Acquisition Heterogeneity}
  {AAAI 2027, under submission}
  {Second Author}
  {TopoCMamba regulates Mamba state propagation with frequency-induced topological conductance, improving anatomical boundary preservation and unseen-domain segmentation across heterogeneous clinical acquisition settings.}
  {
    \item Frequency-Topology Conductance Field for structural and boundary cues.
    \item Conductance-Gated State Scan to suppress cross-boundary state mixing.
    \item Evaluated across 15 datasets and six medical imaging modalities.
  }

\paper
  {LynciaPsyEvo: Enabling Inter-Session Self-Evolution for Personalized LLM Counseling}
  {ICLR 2027, planned submission}
  {First Author}
  {A multi-session counseling agent that keeps the base LLM frozen while adapting therapeutic strategies for each client across sessions, improving personalization and longitudinal self-evolution.}
  {
    \item Per-client counseling strategy adaptation.
    \item PsychEval average score: 7.74, compared with 7.35 for PsychAgent 32B.
  }

\section{Selected Awards and Challenges}
\begin{itemize}
  \item 3rd place, Binary Spike Detection Track, NeuroMM-2026 Challenge.
  \item 3rd place, Five-Class Subtype Classification Track, NeuroMM-2026 Challenge.
\end{itemize}

\section{Research Keywords}
\begin{itemize}
  \item Personalized LLM agents, counseling agents, preference optimization.
  \item EEG analysis, multimodal EEG-video understanding, epileptiform-discharge detection.
  \item Test-time adaptation, domain shift, cross-subject generalization.
  \item Medical image segmentation, Mamba-based architectures, clinical acquisition heterogeneity.
\end{itemize}

\end{document}
